\documentclass{article}
\usepackage{amsmath}
\usepackage{xcolor}
\begin{document}

CAP 04- ECUACIONES DIFERENCIALES DE ORDEN SUPERIOR



E.D. Lineales


\begin{gather*}
\text{\colorbox[RGB]{245,166,35}{$a$}}\text{\colorbox[RGB]{245,166,35}{$_{n}$}}\text{\colorbox[RGB]{245,166,35}{$($}}\text{\colorbox[RGB]{245,166,35}{$x$}}\text{\colorbox[RGB]{245,166,35}{$)$}}y^{n}+\text{\colorbox[RGB]{245,166,35}{$a$}}\text{\colorbox[RGB]{245,166,35}{$_{n-1}$}}\text{\colorbox[RGB]{245,166,35}{$($}}\text{\colorbox[RGB]{245,166,35}{$x$}}\text{\colorbox[RGB]{245,166,35}{$)$}}y^{n-1}+.....+a_{1}(x)y'+\text{\colorbox[RGB]{245,166,35}{$a$}}\text{\colorbox[RGB]{245,166,35}{$_{0}$}}\text{\colorbox[RGB]{245,166,35}{$($}}\text{\colorbox[RGB]{245,166,35}{$x$}}\text{\colorbox[RGB]{245,166,35}{$)$}}y=g(x)
\end{gather*}


La linealidad esta asociada al formato anterior, donde la idea es tener sueltas 

las diferenciales como la variable dependiente, estas pueden estar acompañadas de

funciones de 'x' $\rightarrow$   (las marcadas con naranja) 

\textbf{Clasificacion} :

De acuerdo al miembro de la izq: Concieficientes Constantes, o Sin Coeficientes Ctte

De acuerdo al miembro de la der: Homogeneas o No Homogeneas

Ejemplos:

E.D con coeficientes cttes 


\begin{gather*}
\text{\colorbox[RGB]{248,231,28}{$5$}}\frac{d^{2}y}{dx^{2}}+\text{\colorbox[RGB]{248,231,28}{$10$}}\frac{dy}{dx}-y = .....
\end{gather*}


E.D sin coeficientes cttes 


\begin{gather*}
\text{\colorbox[RGB]{248,231,28}{$5x$}}\text{\colorbox[RGB]{248,231,28}{$^{2}$}}\frac{d^{2}y}{dx^{2}}+\text{\colorbox[RGB]{248,231,28}{$10e$}}\text{\colorbox[RGB]{248,231,28}{$^{2x}$}}\frac{dy}{dx}-y = .....
\end{gather*}


E.D. Homogeneas:


\begin{gather*}
5\frac{d^{2}y}{dx^{2}}+10\frac{dy}{dx}-y =\text{\colorbox[RGB]{248,231,28}{$ 0$}}
\end{gather*}


E.D. NO Homogeneas:


\begin{gather*}
5\frac{d^{2}y}{dx^{2}}+10\frac{dy}{dx}-y = \text{\colorbox[RGB]{248,231,28}{$x$}}\text{\colorbox[RGB]{248,231,28}{$^{2}$}}\text{\colorbox[RGB]{248,231,28}{$+5x-10$}}
\end{gather*}




\textbf{SOLUCION A UNA E.D. de Orden Superior:}

Toda E.D. de orde n $\rightarrow$  se asume que tendra \colorbox[RGB]{248,231,28}{n soluciones linealmente independientes}

Por decir de forma generica, si se toma la E.D. general


\begin{gather*}
\text{\colorbox[RGB]{245,166,35}{$a$}}\text{\colorbox[RGB]{245,166,35}{$_{n}$}}\text{\colorbox[RGB]{245,166,35}{$($}}\text{\colorbox[RGB]{245,166,35}{$x$}}\text{\colorbox[RGB]{245,166,35}{$)$}}y^{n}+\text{\colorbox[RGB]{245,166,35}{$a$}}\text{\colorbox[RGB]{245,166,35}{$_{n-1}$}}\text{\colorbox[RGB]{245,166,35}{$($}}\text{\colorbox[RGB]{245,166,35}{$x$}}\text{\colorbox[RGB]{245,166,35}{$)$}}y^{n-1}+.....+a_{1}(x)y'+\text{\colorbox[RGB]{245,166,35}{$a$}}\text{\colorbox[RGB]{245,166,35}{$_{0}$}}\text{\colorbox[RGB]{245,166,35}{$($}}\text{\colorbox[RGB]{245,166,35}{$x$}}\text{\colorbox[RGB]{245,166,35}{$)$}}y=0 \\
y_{1} = ..... (x^{2}) \\
y_{2}=..... \\
. \\
. \\
y_{n}= .... \\
\text{Por ende cada solucion, al reemplar en la E.D. debe satisfacer} \\
\text{la misma, y generar un valor de 0 al hacer el reeemplazo respectivo} \\
\text{Se dice que la suma de las soluciones, tambien es una solucion} \\
\text{Cuando la E.D. es Homogenea} \\
y_{h}=y_{1}+y_{2}+y_{3}.....y_{n} \\
\text{si la E.D. fuese no Homogenea, se asume} \\
\text{que existe una solucion particular para E.D.} \\
y_{p} = ... \\
\vspace{0.5em} \\
\text{De forma general la solucion como tal seria:} \\
y = y_{h}+y_{p}
\end{gather*}


INDEPENCIA LINEAL

Se dice que un conjunto de funciones y1,y2, y3, son linealmente independientes, \colorbox[RGB]{248,231,28}{cuando no son posibles absoverse entre si operando algebraicamente}.

Por ejemplo, si 
\begin{gather*}
y_{1}= x  ; y_{2} = 5x ;  y_{3}= -3x
\end{gather*}
 determine si son linealmente independiente




\begin{gather*}
y1+y2+y3 = y1+y2+y3 \\
x +5x-3x = 3x  \rightarrow \text{ estas funciones son linealmente dependientes}
\end{gather*}




Por ejemplo, si 
\begin{gather*}
y_{1}= x  ; y_{2} = x^{2} 
\end{gather*}



\begin{gather*}
y1+y2 = y1+y2 \\
x+x^{2} = x+x^{2}  \rightarrow \text{ son funciones linealmente independientes}
\end{gather*}


WRONSKIANO

Es una matriz que se forma a traves del set de funciones dadas, cada fila posterior  a la primera, 

se la calcula con la derivada respectiva

Por decir, si se nos facilita el siguiente set de funciones  {y1,y2 }




\begin{gather*}
W =\matrix{y_{1}}{y_{2}}{y_{1}'}{y_{2}'}=y_{1}*y_{2}'-y_{2}*y_{1}' = ...... \\
W = \matrix{y_{1}}{y_{2}}{y_{1}'}{y_{2}'}{y_{1}''}{y_{2}''}{y_{3}}{y_{3}'}{y_{3}''} \\
Si el valor encontrada, W <> 0  \rightarrow \text{ se tiene la certeza} \\
\text{que el set de funciones son LINEALMENTE INDEPENDIENTE}
\end{gather*}


Dado el siguiente set de funciones 
\begin{gather*}
y_{1}= x  ; y_{2} = x^{2} 
\end{gather*}
, determine el Wronskiano


\begin{gather*}
W =\matrix{x}{x^{2}}{1}{2x}=x*2x'-x^{2}(1) =x^{2} \\
\text{por ende y}_{1} y y_{2}\text{ son }\text{\colorbox[RGB]{248,231,28}{$\text{LINEALMENTE INDEPENDIENTES}$}}
\end{gather*}


Dado el siguiente set de funciones 
\begin{gather*}
y_{1}= 10x^{2}  ; y_{2} = 5x^{2} 
\end{gather*}
, determine el Wronskiano


\begin{gather*}
W =\matrix{10x^{2}}{5x^{2}}{20x}{10x}=100x^{3}-100x^{3}=0 \\
\text{por ende y}_{1} y y_{2}\text{ son LINEALMENTE DEPENDIENTES}
\end{gather*}


E.D. DE 2DO ORDEN


\begin{gather*}
a_{2}(x)y''+a_{1}(x)y'+a_{0}(x)y=0 \\
\text{Por ende se tiene como solucion} \\
y_{1} = ... \\
y_{2} = ... \\
yh = y1+y2 \\
\vspace{0.5em} \\
\text{Asumamos que una solucion generica para este tipo de E.D. considerando que la E.D} \\
\text{ES con Coeficientes CTTE.} \\
Ay''+By'+Cy=0 \\
\text{Se puede considerar que la solucion generica a esta E.D. es} \\
y = e^{rx}   \rightarrow \text{donde r, es un valor real}(raiz) \\
\text{por ende se deberia cumplir el reemplazo con las derivadas, } \\
\text{de tal manera que obtengamos cero } \\
y = e^{rx} \\
y' = re^{rx} \\
y'' = r^{2}e^{rx} \\
\text{Si reemplazamos en la E.D.} \\
Ar^{2}e^{rx}+Bre^{rx}+Ce^{rx}=0 \\
e^{rx}(Ar^{2}+Br+C) =0 \\
Ar^{2}+Br+C =0 \\
\text{ solucion para r} \rightarrow \text{ se puede determinar de varias formas} \\
r1,r2 = \frac{-b±\sqrt{b^{2}-4ac}}{2a} \\
\text{por ende tenemos dos raices, que nos permiten obtener} \\
y1 = e^{r_{1}x} ; y2= e^{r_{2}x}
\end{gather*}


EJERCICIOS


\begin{gather*}
a) \frac{d^{2}y}{dx^{2}}-3\frac{dy}{dx}+2y = 0 \\
y = e^{rx} \\
e^{rx}(r^{2}-3r+2) =0 \\
(r-2)(r-1) =0 \\
r_{1}=2; r_{2}=1 \\
y1 = e^{2x}; y2 =e^{x} \\
y_{h} = e^{2x}+ e^{x} \\
\text{Corroborando, primeramente con y}_{1} \\
y_{1} = e^{2x} \\
y_{1}' = 2e^{2x} \\
y_{2}' = 4e^{2x} \\
(4e^{2x})-3(2e^{2x})+2(e^{2x}) = 0 \\
\vspace{0.5em} \\
y_{3} = 3.14e^{2x} \\
y_{3}' = 6.28e^{2x} \\
y_{3}'' = 12.56e^{2x} \\
( 12.56e^{2x})-3(6.28e^{2x})+2(3.14e^{2x}) = 0 \\
\vspace{0.5em} \\
y = e^{2x}+e^{x}-e^{x}+ 3.14e^{2x} \\
y = c_{1}y_{1}+c_{2}y_{2} \\
y = c_{1}e^{2x}+c_{2}e^{x}
\end{gather*}
 




\begin{gather*}
b) \frac{d^{2}y}{dx^{2}}-4\frac{dy}{dx}+4y = 0 \\
y = e^{rx} \\
e^{rx}(r^{2}-4r+4) =0 \\
(r-2)(r-2) \\
r_{1} =2; r_{2}=2 \\
y_{1}=e^{2x}; y_{2}=e^{2x} \\
y_{h} = c_{1}e^{2x}+c_{2}e^{2x} \\
yh = e^{2x}(c1+c2)   \rightarrow A = c1+c2 \\
yh = Ae^{2x} \\
\text{bajo este planteamiento, la solcuiones se absorven, porque no son LI} \\
y_{h} = c_{1}e^{2x}+c_{2}e^{2x}\text{\colorbox[RGB]{248,231,28}{$x$}}  \\
\text{Se multiplica por x, para garantizar la Independicia lineal}
\end{gather*}


a) 
\begin{gather*}
 y'''+3y''+3y'+y=0
\end{gather*}



\begin{gather*}
y'''+3y''+3y'+y=0 \\
e^{rx}(r^{3}+3r^{2}+3r+1) = 0 \\
\matrix{1}{3}{}{-1}{1}{2}{3}{1}{}{-2}{-1}{-1}{1}{0}{} \\
(r^{2}+2r+1)(r+1)=0 \\
(r+1)(r+1)(r+1) \\
y_{1} = c_{1}e^{-x} \\
y_{2} = c_{2}xe^{-x} \\
y_{3} = c_{3}x^{2}e^{-x} \\
y_{h} = c_{1}e^{-x}+c_{2}\text{\colorbox[RGB]{248,231,28}{$x$}}e^{-x}+c_{3}\text{\colorbox[RGB]{248,231,28}{$x$}}\text{\colorbox[RGB]{248,231,28}{$^{2}$}}e^{-x}
\end{gather*}


b) 
\begin{gather*}
 y'''-y''+y'-y=0.
\end{gather*}



\begin{gather*}
y'''-y''+y'-y=0   \rightarrow    \\
y(0)=1; y'(0)=2; y''(0)=-1 \\
\vspace{0.5em} \\
e^{rx}(r^{3}-r^{2}+r-1) = 0 \\
\matrix{1}{-1}{}{1}{1}{0}{1}{-1}{}{0}{1}{1}{1}{0}{} \\
(r^{2}+1)(r-1)=0 \\
r^{2}+1=0 \\
r^{2}=-1 = r=\sqrt{-1}  \rightarrow i,-i \\
\text{En el caso de determinar raices imaginarias:} \\
r_{1} = a+bi    \rightarrow  r_{2} =a-bi  \\
y = e^{ax}c_{1}cosbx+e^{ax}c_{2}senbx =e^{ax}(c_{1}cosbx+c_{2}\text{senbx})  \\
y_{1} = c_{1}e^{x} \\
y_{2} =c_{2}e^{ix} = c_{2}cosx \\
y_{3} =c_{3}e^{-ix}= c_{3}senx \\
y_{h} = c_{1}e^{x}+c_{2}cosx+c_{3}senx
\end{gather*}


c) 
\begin{gather*}
 y^{(4)}-16y=0
\end{gather*}



\begin{gather*}
 y^{(4)}-16y=0 \\
e^{rx}(r^{4}-16)=0 \\
(r^{2}^{2}-4^{2})=0   \rightarrow  a =r^{2} \\
(a^{2}-4^{2}) = (a-4)(a+4) \\
\text{reemplando el valor de a} \\
(r^{2}-4)(r^{2}+4)=0 \\
(r^{2}-2^{2})(r^{2}+4)=0 \\
(r-2)(r+2)(r^{2}+4)=0   \rightarrow r^{2}=-4 \\
r_{1}=2; r_{2}=-2 ; r_{3}=2i; r_{4}=-2i \\
y_{h} =c_{1}e^{2x}+c_{2}e^{-2x}+ c_{3}cos2x+c_{4}\text{sen2x}
\end{gather*}




e)
\begin{gather*}
 (D^{2}+4D+13)^{3}y=0.
\end{gather*}
  $\rightarrow$  Operador Diferencial: 
\begin{gather*}
\frac{dy}{dx}
\end{gather*}
 = Dy  $\rightarrow$  
\begin{gather*}
D = \frac{d}{dx}
\end{gather*}


por ejemplo:


\begin{gather*}
y''+16y = 0 \rightarrow  \frac{d^{2}y}{dx^{2}} = \frac{d}{dx}\frac{d}{dx}y = D^{2}y \\
D^{2}y+16y=0 \\
y(D^{2}+16)
\end{gather*}


Resolviendo el ejercicio


\begin{gather*}
(D^{2}+4D+13)^{3}y=0   \rightarrow  -2±3i \\
\text{es una E.D. de 6 orden } \rightarrow \text{ se deben determinar 6 raices} \\
(D^{2}+4D+13)(D^{2}+4D+13)(D^{2}+4D+13)=0 \\
D^{2}+4D+13 =0 \\
\text{completando cudrados} \\
(D^{2}+4D+4)-4 +13 \\
(D^{2}+4D+4)+9 =0 \\
(D+ 2)(D+2) +9 =0 \\
(D+ 2)^{2}+9 =0 \\
(D+ 2)^{2} = -9 \\
\text{aplicando raices a ambos elementos} \\
\sqrt{(D+ 2)^{2}}= \sqrt{-9} \\
D+ 2 =±3i \\
D_{1,2} = -2±3i \\
y_{1,2}=e^{-2x}(c_{1}cos3x+c_{2}\text{sen3x}) \\
y_{3,4}=e^{-2x}(c_{3}cos3x+c_{4}\text{sen3x})\text{\colorbox[RGB]{248,231,28}{$x$}} \\
y_{5,6}=e^{-2x}(c_{5}cos3x+c_{6}\text{sen3x})\text{\colorbox[RGB]{248,231,28}{$x$}}\text{\colorbox[RGB]{248,231,28}{$^{2}$}} \\
y_{h} =y_{1,2}+ y_{3,4}+y_{5,6}
\end{gather*}


f) Considerando que una de las soluciones a un E.D. de 6 orden 

esta dada por: 
\begin{gather*}
e^{3x}x^{2}\text{sen2x}
\end{gather*}


1) Determine el resto de las soluciones

2) Determine la E.D. como tal



g) Considerando que una de las soluciones a una ED. de 5to orden esta 

dada por 
\begin{gather*}
x^{2}\text{cos3x}
\end{gather*}
. Es posible determinar la e.d? JUSTIFIQUE su respuesta



h) Considerando que una de las soluciones a una ED. de 5to orden esta 

dada por 
\begin{gather*}
x\text{cos3x}
\end{gather*}
 y adiconalmente se sabe una de las soluciones es

simplemente una contante. Es posible determinar la e.d? JUSTIFIQUE su respuesta

en caso de ser posible, determine el resto de soluciones y la E.D. como tal


\end{document}
